\section*{Abstract}

Traffic monitoring observes road conditions; short-term forecasting anticipates
congestion. This project develops a modular prototype that evaluates the two
tasks on their respective data domains. A vehicle-recognition module fine-tunes
YOLO26m on MIO-TCD and 320 LTA images with teacher pseudo-labels, and is
evaluated on a leakage-aware MIO-TCD split, on LTA teacher agreement, and on a
63-image human-corrected subset of the week-2 validation frames. A per-camera
filter removes fixed false positives, and a stacked ensemble of detector counts,
background differencing and zero-shot CLIP assigns three congestion levels. A
forecasting module compares eight models on PEMSD8 under native and unified
nine-channel inputs, with five-seed checks for the leading four. A workbench
presents both source-labeled outputs.

On MIO-TCD, fine-tuning raises mAP@0.5 from 0.430 to 0.720. On the
human-corrected LTA subset, joint training lowers counting error and raises
recall but does not establish an mAP gain (0.478 versus 0.467). On PEMSD8, GWNET
has the lowest MAE; MultiHeadSTGCN is competitive but not superior. The filter
raises precision from 0.816 to 0.833 at a small recall cost; congestion
estimates agree with 74.3\% of 400 held-out labels. The prototype does not close
the loop from LTA images to local forecasting; calibrated time-aligned series
are still required.
